\documentclass[track=junior,level=3]{kidscode}
\begin{document}
\kctitle{Notes for Families}{What we are doing, how we keep children safe, and how you can help at home}

\section{Dear families,}
Welcome to Bridge to Python, the last level of the junior coding track! Once a week, for 60 minutes,
your child will code in a small group of up to four children (about 8--9 years old, US Grade 3).
Your child has already made animations and games in \textbf{Scratch} in Junior Levels 1 and 2.
This level builds a bridge from those colored blocks to \textbf{Python}, a real programming language
that people use at work every day.

In ten weeks children make art with patterns inside patterns, plan programs in parts, collect data
and draw bar graphs, find out how computers learn from examples, and then write their first Python
programs, in which a little turtle draws on the screen.
\begin{itemize}[itemsep=2pt]
  \item \textbf{Weeks 1--4 (Scratch):} pattern art, planning on paper, a family survey with a bar graph,
        and a game about how computers learn (and how they can learn unfair mistakes).
  \item \textbf{Week 5:} a play day, where children match each Scratch block to its Python twin.
  \item \textbf{Weeks 6--8 (Python):} the free \textbf{Thonny} app, the turtle, repeat loops in Python,
        and their own colors and patterns.
  \item \textbf{Weeks 9--10:} the project, \enquote{My Big Build} (a Scratch game or a Python postcard),
        and a \textbf{graduation showcase} where your child shows it to the family and gets a junior
        graduation certificate.
\end{itemize}
Every session has the same five parts: a warm-up game or short video, a short demonstration
(\enquote{Show me}), building in pairs (\enquote{Make it}), a wiggle break, and a Share where each
child shows what they made. Children work as a pair: the \textbf{Driver} holds the mouse and types,
and the \textbf{Helper} points and reads. In the Python weeks the Helper reads the code aloud letter
by letter, because Python is picky about spelling. They swap when the timer runs out.

\section{Our safety and privacy promises}
\begin{safety}[What we promise]
\begin{itemize}[itemsep=2pt]
  \item \textbf{Offline.} We use the Scratch and Thonny apps on the computer, not websites. Nothing needs the internet.
  \item \textbf{No accounts.} Nobody signs up, logs in or gives an email address. Thonny and Scratch are free.
  \item \textbf{Nicknames only.} Projects and file names show a nickname or a first name, never a full
        name, address, school, age or photo.
  \item \textbf{Surveys ask only what they need.} In the data week children ask family members about
        things like a favorite fruit, never about personal information.
  \item \textbf{AI safety.} Children learn that a computer that learned from examples can be wrong, and
        that they never tell a chatbot or an app their name, address or school.
  \item \textbf{Nothing is uploaded.} Projects stay on the computer. The showcase is at home, with
        family only, on the screen.
  \item \textbf{Healthy screen time.} Every session has time away from the screen and a wiggle break.
  \item \textbf{For everyone.} All pages use color-blind-safe colors with words and icons, and every
        page comes in a dyslexia-friendly edition. Please ask.
\end{itemize}
\end{safety}

\section{What your child needs}
\begin{itemize}[itemsep=2pt]
  \item \textbf{A Mac with Scratch and Thonny.} Both are free. Scratch 3 is in the Mac App Store;
        Thonny comes from \texttt{thonny.org}. No account is needed for either.
  \item \textbf{A mouse.} Dragging blocks with a trackpad is hard for small hands.
  \item \textbf{60 minutes, once a week,} at a time when your child is not too tired or hungry.
  \item \textbf{A grown-up nearby,} to read the pages aloud and to be the Helper when there is one child.
\end{itemize}

\section{How to help at home}
You do not need to know anything about coding to help, not even Python.
\begin{itemize}[itemsep=3pt]
  \item \textbf{Play the play-at-home card.} Each week there is an optional 5--10 minute game, mostly
        without a computer. Stop while it is still fun.
  \item \textbf{Be a survey answer.} In week 3 your child may ask you a survey question. Answer, and
        ask what the bar graph shows.
  \item \textbf{Red text is not a disaster.} When Python finds a mistake, it shows red text. Ask
        \enquote{What does the last red line say?} and look for the spelling slip together.
  \item \textbf{Praise effort.} \enquote{You kept trying until it worked!} means more than
        \enquote{You're so clever.}
  \item \textbf{Hands off the keyboard.} If you help, be the Helper: point and spell, and let your child type.
\end{itemize}

\section{Ask your child about\dots}
One question for each week. They make great conversation starters at dinner.

\begin{center}\small\renewcommand{\arraystretch}{1.45}
\begin{tabularx}{\linewidth}{@{}>{\bfseries}p{1.0cm}>{\RaggedRight}p{4.2cm}>{\RaggedRight\arraybackslash}X@{}}
\toprule
Week & Topic & Ask your child\dots\\
\midrule
1 & Patterns in Patterns & \enquote{Can you clap a pattern inside a pattern? How did the cat draw a flower of squares?}\\
2 & Plan It on Paper & \enquote{How would you plan a party in parts? Why test each part?}\\
3 & Data Detectives & \enquote{What did your survey find out? Which bar was the tallest?}\\
4 & How Do Computers Learn? & \enquote{How did you teach the robot? Why did it make a mistake?}\\
5 & Play Day: Blocks to Words & \enquote{What is the Python twin of the move block?}\\
6 & Hello, Python! & \enquote{What did your turtle draw? What was the funniest typo you found?}\\
7 & Repeat in Python & \enquote{How do you say \enquote{repeat 4} in Python? Why does the indent matter?}\\
8 & Colors and My Own Pattern & \enquote{Which colors and numbers did you change? Whose idea did you remix?}\\
9 & Project Studio & \enquote{Did you choose Scratch or Python for your Big Build? What is on your plan card?}\\
10 & Graduation Showcase & \enquote{What are you proudest of? What would you like to code next?}\\
\bottomrule
\end{tabularx}
\end{center}

\section{How to celebrate}
\begin{itemize}[itemsep=3pt]
  \item \textbf{Come to the graduation showcase} in week 10. Your child will show the Big Build, tell
        you what one part of the code does and maybe show you a bug they fixed. You get sticker feedback
        cards to fill in.
  \item \textbf{Graduation certificates.} Every child who finishes Bridge to Python gets a printed junior
        graduation certificate.
  \item \textbf{A fridge gallery.} Ask your child to draw their turtle picture or game for the fridge.
        Any photos or screenshots stay at home, not online.
  \item \textbf{What comes next.} After the junior track, children are ready for the main course's
        Level 1, which continues in Scratch and Python.
\end{itemize}

\bigskip
Thank you for coding with us. If anything would help your child learn (for example bigger text,
more breaks or a different time of day), please tell me.

\bigskip
With warm wishes,\par\bigskip
\rule{6cm}{0.4pt}\par
{\small (Your coding club host)}
\end{document}
